% !TEX program = pdflatex
\documentclass[10pt,aspectratio=169]{beamer}

\usepackage[utf8]{inputenc}
\usepackage[T1]{fontenc}
\usepackage[english]{babel}
\usepackage{lmodern}
\usepackage{microtype}
\usepackage{amsmath,amssymb,bm,mathtools}
\usepackage{booktabs,array}

\usetheme[numbering=fraction,progressbar=frametitle]{metropolis}
\setbeamertemplate{frame numbering}[fraction]
\setbeamertemplate{navigation symbols}{}

\definecolor{CourseBlue}{HTML}{1A3C68}
\definecolor{CourseGreen}{HTML}{216E39}
\definecolor{CourseGold}{HTML}{B8860B}
\definecolor{CourseRed}{HTML}{9A2C2C}
\definecolor{SoftBlue}{HTML}{F2F6FA}
\definecolor{SoftGreen}{HTML}{F1F8F3}
\definecolor{SoftRed}{HTML}{FFF3F3}

\setbeamercolor{normal text}{fg=black,bg=white}
\setbeamercolor{structure}{fg=CourseBlue}
\setbeamercolor{title}{fg=CourseBlue}
\setbeamerfont{title}{size=\Large,series=\bfseries}
\setbeamercolor{frametitle}{bg=CourseBlue,fg=white}
\setbeamerfont{frametitle}{series=\bfseries,size=\large}
\setbeamertemplate{frametitle}{
  \nointerlineskip
  \begin{beamercolorbox}[wd=\paperwidth,ht=3.2ex,dp=1.4ex,leftskip=1em,rightskip=1em]{frametitle}
    \usebeamerfont{frametitle}\insertframetitle
  \end{beamercolorbox}
}
\setbeamercolor{block title}{bg=CourseBlue,fg=white}
\setbeamercolor{block body}{bg=SoftBlue,fg=black}
\setbeamercolor{block title example}{bg=CourseGreen,fg=white}
\setbeamercolor{block body example}{bg=SoftGreen,fg=black}
\setbeamercolor{block title alerted}{bg=CourseRed,fg=white}
\setbeamercolor{block body alerted}{bg=SoftRed,fg=black}
\setbeamertemplate{blocks}[rounded][shadow=false]

\newenvironment{modelbox}{\begin{beamercolorbox}[rounded=true,sep=0.7em,wd=\textwidth]{block body}
  \textbf{\color{CourseBlue}MODEL}\quad}{\end{beamercolorbox}}
\newenvironment{marketbox}{\begin{beamercolorbox}[rounded=true,sep=0.7em,wd=\textwidth]{block body example}
  \textbf{\color{CourseGreen}NO-ARBITRAGE}\quad}{\end{beamercolorbox}}
\newenvironment{simplificationbox}{\begin{beamercolorbox}[rounded=true,sep=0.7em,wd=\textwidth]{block body alerted}
  \textbf{\color{CourseRed}WATCH OUT}\quad}{\end{beamercolorbox}}
\newenvironment{pausebox}{\begin{beamercolorbox}[rounded=true,sep=0.8em,wd=\textwidth]{block body}
  \textbf{\color{CourseGold}YOUR TURN}\quad}{\end{beamercolorbox}}

\newcommand{\E}{\mathbb E}
\newcommand{\Q}{\mathbb Q}
\newcommand{\F}{\mathcal F}
\newcommand{\dd}{\,\mathrm d}
\newcommand{\Ncdf}{\Phi}

\AtBeginSection[]{
  \begin{frame}[plain]\vfill
    {\Large\bfseries\color{CourseBlue}\insertsectionhead\par}
    \medskip{\color{CourseBlue}\rule{0.92\textwidth}{0.8pt}}\vfill
  \end{frame}
}

\title{Beyond Black--Scholes}
\subtitle{Local volatility, stochastic volatility, and stochastic rates}
\author{}
\date{Lecture 12}

\begin{document}

\begin{frame}[plain]
  \vfill
  {\LARGE\bfseries\color{CourseBlue}Beyond Black--Scholes\par}
  \medskip{\large Local volatility, stochastic volatility, and stochastic rates\par}
  \medskip{\color{CourseGold}\rule{0.94\textwidth}{0.8pt}}
  \medskip{Lecture 12\par}\vfill
\end{frame}

\begin{frame}{The model must explain the prices we observe}
\begin{block}{The question}
Black--Scholes is a powerful benchmark. Which of its assumptions fail in the
market, and what model extension addresses each failure?
\end{block}
\begin{enumerate}
  \item Read the implied-volatility surface as evidence against constant
  volatility.
  \item Build local volatility to fit today's vanilla-option surface.
  \item Make volatility random to describe its future dynamics.
  \item Make discounting random when the yield curve is a source of risk.
\end{enumerate}
\begin{marketbox}
The central trade-off is always the same: fit, realism, hedging, and
tractability cannot all be maximized at once.
\end{marketbox}
\end{frame}

\section{Why Black--Scholes is not enough}

\begin{frame}{One constant volatility cannot fit a whole option market}
Black--Scholes assumes, under \(\Q\),
\[
\frac{\dd S_t}{S_t}=r\dd t+\sigma\dd W_t,
\qquad \sigma\text{ constant}.
\]
If the model were exact, Black--Scholes implied volatility would be the same
for every strike \(K\) and maturity \(T\).
\begin{columns}[T,onlytextwidth]
\column{0.48\textwidth}
\begin{block}{What markets show}
Implied volatility varies by strike: smiles and, in equities, a strong
downside skew.
\end{block}
\column{0.48\textwidth}
\begin{exampleblock}{What markets also show}
Implied volatility varies by maturity; uncertainty has its own term
structure.
\end{exampleblock}
\end{columns}
\begin{simplificationbox}
A single number \(\sigma\) cannot reproduce a two-dimensional surface
\((K,T)\mapsto\sigma_{\mathrm{imp}}(K,T)\).
\end{simplificationbox}
\end{frame}

\begin{frame}{The equity skew has an economic interpretation}
For a fixed maturity, a typical equity surface has
\[
\sigma_{\mathrm{imp}}(K_{\mathrm{low}})
>
\sigma_{\mathrm{imp}}(K_{\mathrm{ATM}})
>
\sigma_{\mathrm{imp}}(K_{\mathrm{high}}).
\]
\begin{block}{Why low-strike puts are expensive}
They protect against severe equity declines. Strong demand for this insurance
raises their market prices and therefore their implied volatilities.
\end{block}
\begin{pausebox}
Does a high implied volatility at a low strike mean that the historical
volatility parameter literally changed for that option?
\end{pausebox}
\pause
\begin{exampleblock}{Answer}
No. It is the Black--Scholes parameter that matches a market price.
\end{exampleblock}
\end{frame}

\begin{frame}{Three extensions answer three different questions}
\begin{center}
\begin{tabular}{p{0.27\textwidth}|p{0.29\textwidth}|p{0.29\textwidth}}
\toprule
Observed issue & Model extension & Main objective\\
\midrule
Static smile across \(K,T\) & Local volatility & Fit vanilla prices today\\
Volatility moves through time & Stochastic volatility & Model smile dynamics and volatility risk\\
Discounting and curve move & Stochastic rates & Model bonds and the term structure\\
\bottomrule
\end{tabular}
\end{center}
\begin{marketbox}
These models do not replace no-arbitrage pricing. They replace restrictive
assumptions about what is deterministic or constant.
\end{marketbox}
\end{frame}

\section{Local volatility fits today's surface}

\begin{frame}{Local volatility makes the diffusion coefficient state dependent}
Under the risk-neutral measure,
\[
\boxed{
\frac{\dd S_t}{S_t}=r\dd t+
\sigma_{\mathrm{loc}}(t,S_t)\dd W_t.}
\]
\begin{columns}[T,onlytextwidth]
\column{0.48\textwidth}
\begin{block}{What remains}
The model is one-factor and Markovian. Given \(t\) and \(S_t\), the future
law is determined.
\end{block}
\column{0.48\textwidth}
\begin{exampleblock}{What changes}
Instantaneous variance is no longer constant:
\[
\operatorname{Var}\!\left(\frac{\dd S_t}{S_t}\middle|\F_t\right)
=\sigma_{\mathrm{loc}}(t,S_t)^2\dd t.
\]
\end{exampleblock}
\end{columns}
\begin{marketbox}
Black--Scholes is the special case
\(\sigma_{\mathrm{loc}}(t,s)\equiv\sigma\).
\end{marketbox}
\end{frame}

\begin{frame}{The pricing PDE changes in one precise place}
For a European payoff \(g(S_T)\), the price \(u(t,s)\) solves
\[
u_t+rs\,u_s+\tfrac12\sigma_{\mathrm{loc}}(t,s)^2s^2u_{ss}-ru=0,
\qquad u(T,s)=g(s).
\]
\begin{block}{Compare with Black--Scholes}
The pricing logic, risk-neutral drift, and terminal condition are unchanged.
Only the constant variance \(\sigma^2\) becomes the local quantity
\(\sigma_{\mathrm{loc}}(t,s)^2\).
\end{block}
\begin{pausebox}
Does local volatility introduce a new random state variable in addition to
the stock price?
\end{pausebox}
\pause
\begin{exampleblock}{Answer}
No. It is a deterministic function of \(t,S_t\), so the model remains
one-dimensional.
\end{exampleblock}
\end{frame}

\begin{frame}{Dupire extracts local volatility from call prices}
Let \(C(T,K)\) be the market price of a call by maturity and strike. With
constant rates and no dividends,
\[
\boxed{
\sigma_{\mathrm{loc}}(T,K)^2
=
\frac{C_T(T,K)+rK\,C_K(T,K)}
{\tfrac12K^2C_{KK}(T,K)}.}
\]
\begin{columns}[T,onlytextwidth]
\column{0.48\textwidth}
\begin{block}{Denominator}
\[
C_{KK}(T,K)
\]
is the curvature in strike; it is linked to the risk-neutral density.
\end{block}
\column{0.48\textwidth}
\begin{exampleblock}{Numerator}
\[
C_T+rKC_K
\]
captures how the call surface changes with maturity and strike.
\end{exampleblock}
\end{columns}
\end{frame}

\begin{frame}{Call convexity is a no-arbitrage input to Dupire}
For an arbitrage-free call-price surface,
\[
C_{KK}(T,K)\geq0.
\]
Formally,
\[
C_{KK}(T,K)=e^{-rT}p_{S_T}^{\Q}(K),
\]
where \(p_{S_T}^{\Q}\) is the risk-neutral density.
\begin{pausebox}
Why is applying Dupire's formula directly to a sparse table of noisy market
quotes dangerous?
\end{pausebox}
\pause
\begin{alertblock}{Answer}
The formula differentiates the surface, especially twice in strike.
Interpolation noise is amplified and can produce negative or unstable local
variances unless the input surface is smoothed and arbitrage controlled.
\end{alertblock}
\end{frame}

\begin{frame}{Local volatility calibrates exactly, but only statically}
A practical calibration pipeline is:
\[
\text{quotes}\longrightarrow
\text{arbitrage-free call or implied-vol surface}
\longrightarrow
\text{smooth derivatives}
\longrightarrow
\sigma_{\mathrm{loc}}(T,K).
\]
\begin{columns}[T,onlytextwidth]
\column{0.48\textwidth}
\begin{exampleblock}{Strength}
With a clean surface, local volatility can reproduce today’s European
vanilla prices across the quoted grid.
\end{exampleblock}
\column{0.48\textwidth}
\begin{alertblock}{Limitation}
It prescribes future smile motion through a deterministic function of spot
and time; those dynamics are often not realistic.
\end{alertblock}
\end{columns}
\end{frame}

\section{Stochastic volatility models volatility itself}

\begin{frame}{A random variance factor creates new dynamics}
A generic stochastic-volatility model has
\[
\frac{\dd S_t}{S_t}=r\dd t+\sqrt{V_t}\dd W_t^{(1)},
\]
\[
\dd V_t=b(t,V_t)\dd t+a(t,V_t)\dd W_t^{(2)},
\qquad
\dd\langle W^{(1)},W^{(2)}\rangle_t=\rho\dd t.
\]
\begin{block}{The state is now \((S_t,V_t)\)}
Even if two paths have the same spot \(S_t\), their option values can differ
when their current variance levels \(V_t\) differ.
\end{block}
\begin{marketbox}
The pair \((S,V)\) is Markovian, but the stock alone is generally not.
\end{marketbox}
\end{frame}

\begin{frame}{Heston adds mean reversion, vol-of-vol, and correlation}
A standard Heston specification is
\[
\frac{\dd S_t}{S_t}=r\dd t+\sqrt{V_t}\dd W_t^{(1)},
\qquad
\dd V_t=\kappa(\theta-V_t)\dd t+\eta\sqrt{V_t}\dd W_t^{(2)}.
\]
\begin{center}
\begin{tabular}{c|l}
\toprule
Parameter & Economic role\\
\midrule
\(\theta\) & long-run variance level\\
\(\kappa\) & speed of mean reversion toward \(\theta\)\\
\(\eta\) & volatility of volatility\\
\(\rho\) & correlation between spot and variance shocks\\
\bottomrule
\end{tabular}
\end{center}
\end{frame}

\begin{frame}{Negative spot-volatility correlation generates an equity skew}
\begin{pausebox}
In an equity market, a negative return shock often arrives with a higher
volatility level. What sign of \(\rho\) captures this leverage effect?
\end{pausebox}
\pause
\begin{exampleblock}{Answer}
\[
\boxed{\rho<0.}
\]
A negative stock shock tends to be accompanied by a positive variance shock.
Low-strike puts become relatively valuable, generating a downward implied-
volatility skew.
\end{exampleblock}
\begin{marketbox}
Random variance creates smile curvature; correlation controls much of its
asymmetry.
\end{marketbox}
\end{frame}

\begin{frame}{Mean reversion gives a useful variance forecast}
From
\[
\dd V_t=\kappa(\theta-V_t)\dd t+\eta\sqrt{V_t}\dd W_t,
\]
the conditional mean is
\[
\boxed{
\E[V_T\mid V_t=v]=\theta+(v-\theta)e^{-\kappa(T-t)}.}
\]
\begin{pausebox}
Suppose \(v>\theta\). Is the expected future variance above or below today’s
variance, and why?
\end{pausebox}
\pause
\begin{exampleblock}{Answer}
Below \(v\), because the drift \(\kappa(\theta-v)\) is negative. The larger
\(\kappa\), the faster the expected pull toward the long-run level.
\end{exampleblock}
\end{frame}

\begin{frame}{More realism makes hedging and pricing harder}
\begin{columns}[T,onlytextwidth]
\column{0.48\textwidth}
\begin{block}{Pricing}
The option value is \(u(t,s,v)\), so the PDE has two state variables.
Characteristic-function methods, PDEs, and Monte Carlo are commonly used.
\end{block}
\column{0.48\textwidth}
\begin{alertblock}{Hedging}
There are two sources of risk but often only the stock and cash are traded.
The model is generally incomplete; variance risk cannot be perfectly hedged
with the stock alone.
\end{alertblock}
\end{columns}
\begin{marketbox}
In practice, liquid vanilla options are often used alongside the stock to
manage volatility exposure.
\end{marketbox}
\end{frame}

\begin{frame}{Local and stochastic volatility solve different problems}
\begin{center}
\begin{tabular}{p{0.30\textwidth}|p{0.29\textwidth}|p{0.27\textwidth}}
\toprule
 & Local volatility & Stochastic volatility\\
\midrule
Volatility & deterministic \(\sigma_{\mathrm{loc}}(t,S)\) & random factor \(V_t\)\\
Main strength & exact static vanilla fit & realistic smile dynamics\\
Dimension & one factor & at least two factors\\
Hedging & closer to complete diffusion hedging & generally incomplete\\
\bottomrule
\end{tabular}
\end{center}
\begin{pausebox}
Which model is more naturally suited to a product sensitive to future
realized volatility rather than only today’s vanilla prices?
\end{pausebox}
\pause
\begin{exampleblock}{Answer}
Stochastic volatility: it retains a separate random variance state and is
designed to describe its future evolution.
\end{exampleblock}
\end{frame}

\section{Stochastic rates make discounting risky}

\begin{frame}{With stochastic rates, the discount factor is random}
If the short rate is \(r_t\), the pathwise discount factor is
\[
D(t,T)=\exp\left(-\int_t^T r_s\dd s\right).
\]
The zero-coupon bond is its risk-neutral value:
\[
\boxed{
P(t,T)=\E^\Q\!\left[D(t,T)\mid\F_t\right].}
\]
\begin{block}{Key distinction}
With constant rates, \(P(t,T)=e^{-r(T-t)}\) is deterministic. With
stochastic rates, discounting itself is a source of uncertainty.
\end{block}
\end{frame}

\begin{frame}{Bond prices, zero rates, and forwards encode the same curve}
For maturity \(T>t\),
\[
P(t,T)=e^{-R(t,T)(T-t)},
\qquad
R(t,T)=-\frac{\log P(t,T)}{T-t},
\]
and the instantaneous forward rate is
\[
f(t,T)=-\partial_T\log P(t,T).
\]
\begin{marketbox}
The bond-price curve \(T\mapsto P(t,T)\), zero-rate curve \(R(t,T)\), and
forward curve \(f(t,T)\) are three parameterizations of the same term
structure.
\end{marketbox}
\end{frame}

\begin{frame}{Worked curve example: extract a zero rate and a forward rate}
Suppose
\[
P(0,1)=0.96,\qquad P(0,2)=0.90.
\]
The two-year continuously compounded zero rate is
\[
R(0,2)=-\frac12\log(0.90)=5.27\%.
\]
The simply compounded one-year forward rate for \([1,2]\) is
\[
L(0;1,2)=\frac{P(0,1)/P(0,2)-1}{2-1}
=\frac{0.96}{0.90}-1=6.67\%.
\]
\begin{block}{Reading the numbers}
The forward rate exceeds the two-year average zero rate: the curve embeds
higher implied rates over the second year than over the first.
\end{block}
\end{frame}

\begin{frame}{Coupon bonds are portfolios of zero-coupon bonds}
A bond paying coupons \(c_i\) at dates \(T_i\), and principal \(F\) at
\(T_n\), is worth
\[
\boxed{
B_t=\sum_{i=1}^n c_iP(t,T_i)+F P(t,T_n).}
\]
\begin{pausebox}
A two-year bond pays \(4\) at year 1 and \(104\) at year 2. If
\(P(0,1)=0.96\) and \(P(0,2)=0.90\), what is its price?
\end{pausebox}
\pause
\begin{exampleblock}{Answer}
\[
B_0=4(0.96)+104(0.90)=97.44.
\]
The entire pricing problem reduces to discount factors at the cash-flow
dates.
\end{exampleblock}
\end{frame}

\begin{frame}{Short-rate models compress the curve into a few factors}
A short-rate model starts from a process \(r_t\), then derives bond prices:
\[
P(t,T)=\E^\Q\!\left[
\exp\left(-\int_t^T r_s\dd s\right)\middle|\F_t\right].
\]
\begin{columns}[T,onlytextwidth]
\column{0.48\textwidth}
\begin{block}{Vasicek}
\[
\dd r_t=a(b-r_t)\dd t+\sigma\dd W_t.
\]
Gaussian and mean reverting, but rates can become negative.
\end{block}
\column{0.48\textwidth}
\begin{exampleblock}{CIR}
\[
\dd r_t=a(b-r_t)\dd t+\sigma\sqrt{r_t}\dd W_t.
\]
Mean reverting with a square-root diffusion that supports positivity.
\end{exampleblock}
\end{columns}
\end{frame}

\begin{frame}{HJM models the forward curve directly}
Instead of modeling only \(r_t\), Heath--Jarrow--Morton specifies
\[
\dd f(t,T)=\alpha(t,T)\dd t+\sigma(t,T)\dd W_t.
\]
\begin{block}{The no-arbitrage restriction}
Under the risk-neutral measure, the drift \(\alpha(t,T)\) is not freely
chosen. Once forward-rate volatility is specified, no-arbitrage determines
the drift.
\end{block}
\begin{columns}[T,onlytextwidth]
\column{0.48\textwidth}
\begin{exampleblock}{Strength}
It starts from the object the market quotes: the whole forward curve.
\end{exampleblock}
\column{0.48\textwidth}
\begin{alertblock}{Cost}
Without additional structure, the model is infinite-dimensional.
\end{alertblock}
\end{columns}
\end{frame}

\begin{frame}{When must rates be modelled stochastically?}
\begin{center}
\begin{tabular}{p{0.39\textwidth}|p{0.40\textwidth}}
\toprule
Often a reasonable approximation & Stochastic-rate model matters\\
\midrule
Short-dated vanilla equity option & Bond or bond option\\
Small rate sensitivity & Caplet, floorlet, or swaption\\
Short horizon & Long-dated equity or hybrid payoff\\
Deterministic discounting acceptable & Discounting is itself a key risk\\
\bottomrule
\end{tabular}
\end{center}
\begin{pausebox}
Which column fits a five-year callable bond: left or right?
\end{pausebox}
\pause
\begin{exampleblock}{Answer}
Right. Its cash flows and exercise decision are directly exposed to the
evolution of the yield curve.
\end{exampleblock}
\end{frame}

\section{Calibration and model choice}

\begin{frame}{Calibration targets depend on the model class}
\begin{center}
\begin{tabular}{p{0.27\textwidth}|p{0.29\textwidth}|p{0.29\textwidth}}
\toprule
Model class & What is fitted & Typical challenge\\
\midrule
Local volatility & vanilla implied-vol surface & smooth, arbitrage-free derivatives\\
Stochastic volatility & surface through \(\kappa,\theta,\eta,\rho,V_0\) & non-linear optimization and stability\\
Stochastic rates & initial curve plus caps, swaptions, bond options & fit curve and dynamics together\\
\bottomrule
\end{tabular}
\end{center}
\begin{marketbox}
Calibration is not merely parameter estimation. It aligns a model with
tradable market prices before the model is used for valuation or risk.
\end{marketbox}
\end{frame}

\begin{frame}{Model choice is a product decision, not a beauty contest}
\begin{pausebox}
You must value a long-dated equity-linked note whose payoff depends on both
the stock and future discount factors. Is local volatility alone enough?
\end{pausebox}
\pause
\begin{exampleblock}{Answer}
Usually not. The product is exposed both to volatility dynamics and to rate
risk. A hybrid framework, or at least explicit stochastic-rate modelling,
may be necessary.
\end{exampleblock}
\begin{block}{Practical rule}
Use the least complex model that captures the material risk factors and can
be calibrated reliably to liquid hedging instruments.
\end{block}
\end{frame}

\begin{frame}{The three extensions form one modeling hierarchy}
\[
\boxed{
\text{market prices}
\longrightarrow
\text{model assumptions}
\longrightarrow
\text{calibration}
\longrightarrow
\text{pricing and hedging}.}
\]
\begin{enumerate}
  \item Local volatility fits the current vanilla surface.
  \item Stochastic volatility adds a random volatility state and richer
  dynamics.
  \item Stochastic rates add a random discount curve and fixed-income risk.
  \item More realism comes with more state variables, calibration work, and
  hedging uncertainty.
\end{enumerate}
\end{frame}

\begin{frame}{Final check: start from the risk you need to model}
\begin{pausebox}
A desk is hedging a barrier option and needs to reproduce today’s full
vanilla smile surface accurately. Which model is the natural first
benchmark, and what caveat should remain in mind?
\end{pausebox}
\pause
\begin{exampleblock}{Answer}
Local volatility is the natural first benchmark because it fits today's
vanilla surface. Its future smile dynamics may nevertheless be unrealistic,
so the hedge can be model dependent.
\end{exampleblock}
\vfill
\begin{center}
{\large\color{CourseBlue}\bfseries A richer model earns its place only by explaining a material risk.}
\end{center}
\end{frame}

\end{document}
